\documentclass[11pt,a4paper]{article}
\usepackage[margin=1in]{geometry}
\usepackage{times}
\usepackage{latexsym}
\usepackage{amsmath}
\usepackage{amssymb}
\usepackage{booktabs}
\usepackage{multirow}
\usepackage{graphicx}
\usepackage{xcolor}
\usepackage{enumitem}
\usepackage{natbib}
\usepackage[hidelinks]{hyperref}

\newcommand{\css}{\textsc{CSS}}
\newcommand{\cpr}{\textsc{CPR}}
\newcommand{\cmd}{\textsc{CMD}}

\title{When Does Cross-Method Consensus Improve Circuit Discovery?\\A Systematic Analysis of Failure Modes}

\author{Anonymous}

\begin{document}
\maketitle

\begin{abstract}
Mechanistic interpretability increasingly relies on automated circuit discovery to identify minimal subgraphs that explain model behavior. When multiple discovery methods exist, a natural hypothesis is that their consensus---edges selected by all methods---yields more faithful circuits. We test this hypothesis systematically by applying binary consensus scoring (\css{}) across three methods (EAP-IG, Edge Pruning, ACDC) on four tasks and two model scales (GPT-2 Small, 32K edges; Qwen-0.5B, 179K edges). Consensus produces a marginal advantage on only 1 of 5 configurations (IOI on GPT-2, $\Delta\cpr{}=0.08$, $p=0.09$), which attenuates to $\Delta=0.02$ ($p=0.27$) after controlling for circuit size. We identify three distinct failure modes: (1) ACDC's \texttt{isinf} binarization collapses consensus to a single method's output, (2) high inter-method agreement (Jaccard $>0.5$) renders consensus redundant, and (3) component-level scoring granularity on larger models produces near-zero cross-method overlap. These findings yield practical diagnostic guidelines: practitioners should check pairwise Jaccard similarity and scoring granularity compatibility \emph{before} investing in multi-method consensus pipelines.
\end{abstract}

%===================================================================
\section{Introduction}
%===================================================================

Circuit discovery aims to identify minimal computational subgraphs within neural networks that are sufficient to reproduce a target behavior \cite{conmy2023towards,wang2023interpretability}. Several automated methods now exist---gradient-based approaches such as EAP-IG \cite{syed2023attribution} and Edge Pruning \cite{bhaskar2024finding}, ablation-based methods such as ACDC \cite{conmy2023towards}, and hybrid approaches like Iterative Feature Refinement (IFR).

When multiple methods are available for the same problem, ensemble strategies from machine learning suggest that their agreement should be more reliable than any individual method. This intuition is formalized in feature selection stability theory \cite{nogueira2016measuring}: the intersection of independently selected feature sets should be enriched for genuinely important features. Applied to circuit discovery, this motivates \emph{Cross-method Consensus Scoring} (\css{}): binarize each method's output and take their intersection.

We test this hypothesis systematically across four tasks (Indirect Object Identification, Greater-Than Arithmetic, Subject-Verb Agreement, Country-Capital Factual Recall) and two model scales (GPT-2 Small with ${\sim}32$K edges, Qwen-0.5B with ${\sim}179$K edges). Our main findings are:

\begin{enumerate}[nosep,leftmargin=*]
    \item \css{} produces a marginal faithfulness advantage on only 1 of 5 configurations (IOI on GPT-2), with the advantage largely attributable to circuit size effects (\S\ref{sec:size-confound}).
    \item We identify three mechanistically distinct failure modes: binarization pathology (\S\ref{sec:binarization}), consensus redundancy under high agreement (\S\ref{sec:redundancy}), and scoring granularity mismatch at scale (\S\ref{sec:granularity}).
    \item A simple Jaccard-based diagnostic correctly predicts \css{} utility in 4 of 5 configurations, though statistical power is insufficient for formal validation (\S\ref{sec:diagnostic}).
\end{enumerate}

Our contribution is not the consensus framework itself---which follows directly from stability indices \cite{nogueira2016measuring,kuncheva2003measures}---but the systematic empirical characterization of \emph{when and why} it fails in the circuit discovery setting, yielding actionable practitioner guidelines.

%===================================================================
\section{Background and Related Work}
%===================================================================

\paragraph{Circuit discovery methods.}
We study three representative methods spanning two paradigms. \textbf{EAP-IG} \cite{syed2023attribution} computes integrated-gradient attributions for each edge in the computational graph. \textbf{Edge Pruning} (EP; \citealp{bhaskar2024finding}) optimizes continuous edge masks via gradient descent. Both are gradient-based and score at edge-level granularity. \textbf{ACDC} \cite{conmy2023towards} iteratively removes edges via activation patching, scoring at node level and propagating scores to edges.

\paragraph{Consensus and ensemble methods.}
The idea of aggregating multiple discovery methods has been explored in concurrent work. \citet{mondorf2025cross} propose cross-method score averaging over continuous attribution values. PIE \cite{chen2026pie} uses iterative pruning ensembles within a single paradigm. CIRCUS \cite{parekh2026circus} applies within-method threshold diversity rather than cross-method consensus. Our approach differs in using binary consensus (intersection), which enables direct failure mode analysis without conflating score-scale incompatibilities.

\paragraph{Feature selection stability.}
Binary consensus is an instance of the stability index framework of \citet{nogueira2016measuring}. The theoretical conditions under which ensemble intersection improves over individual selectors are well-characterized \cite{kuncheva2003measures}: methods must be (a) individually accurate and (b) sufficiently diverse (low pairwise correlation). Our failure modes map directly onto violations of these conditions.

%===================================================================
\section{Method: Cross-Method Consensus Scoring}
\label{sec:method}
%===================================================================

\paragraph{Setup.}
Given a language model $\mathcal{M}$ with computational graph $G=(V,E)$ where $|E|$ can range from ${\sim}32$K (GPT-2 Small) to ${\sim}179$K (Qwen-0.5B), and a set of $K$ circuit discovery methods $\{M_1, \ldots, M_K\}$, each method produces a ranking $r_k: E \to \mathbb{R}$.

\paragraph{Binarization.}
We convert each continuous ranking to a binary mask. For gradient-based methods (EAP-IG, EP), we select the top-$N$ edges by attribution magnitude. For ACDC, which produces a binary circuit, we use its native output (which we show in \S\ref{sec:binarization} to be problematic). As a normalization, we also evaluate top-$N$ binarization uniformly across all methods.

\paragraph{Consensus scoring.}
For each edge $e \in E$, we compute:
\begin{equation}
    \css{}(e) = \frac{1}{K} \sum_{k=1}^{K} \mathbf{1}[e \in \text{top-}N(M_k)]
\end{equation}
Edges with $\css{}=1.0$ (selected by all $K$ methods) form the consensus circuit. We also examine $\css{}\geq 2/3$ (majority consensus).

\paragraph{Evaluation metrics.}
We evaluate circuits using \textbf{Circuit Performance Recovery} (\cpr{}): the fraction of clean model performance recovered when running only the circuit edges (with remaining edges ablated). We use zero ablation throughout (see \S\ref{sec:limitations} for discussion). We also report \textbf{Circuit Mean Degradation} (\cmd{} $= 1 - \cpr{}$).

%===================================================================
\section{Experiments}
\label{sec:experiments}
%===================================================================

We evaluate \css{} with $K=3$ methods (EAP-IG, EP, ACDC) across four tasks on GPT-2 Small, and additionally test scaling behavior on Qwen-0.5B for IOI.

%-------------------------------------------------------------------
\subsection{IOI: The Marginal Positive Case}
\label{sec:ioi}
%-------------------------------------------------------------------

Indirect Object Identification (IOI; \citealp{wang2023interpretability}) is a well-studied task where the model must identify ``Mary'' in ``When John and Mary went to the store, John gave a drink to \underline{\hspace{1em}}.''

On GPT-2 Small with top-100 binarization, the \css{}$=1.0$ consensus circuit contains only 7 edges. In the original (uncontrolled) comparison, this circuit achieves $\cpr{}=0.109$, outperforming all single-method top-100 circuits. However, this comparison conflates the consensus effect with circuit size: a 7-edge circuit naturally suffers less ablation damage than a 100-edge circuit.

\paragraph{Size confound analysis.}
\label{sec:size-confound}
To isolate the consensus effect from circuit size, we compare \css{}$=1.0$ (7 edges) against the top-7 edges from each individual method and a random baseline. Table~\ref{tab:ioi-size} presents the results.

\begin{table}[t]
\centering
\small
\begin{tabular}{lcccc}
\toprule
\textbf{Circuit} & \textbf{N} & \textbf{\cpr{}} & \textbf{\cmd{}} & \textbf{$\Delta$} \\
\midrule
\css{}$=1.0$ & 7 & 0.109 & 0.911 & --- \\
ACDC top-7 & 7 & 0.092 & 0.928 & $-0.017$ \\
EAP-IG top-7 & 7 & 0.078 & 0.942 & $-0.031$ \\
EP top-7 & 7 & 0.065 & 0.955 & $-0.044$ \\
Random-7 & 7 & $-0.031_{\pm 0.045}$ & $1.051_{\pm 0.045}$ & $-0.140$ \\
\bottomrule
\end{tabular}
\caption{IOI size-controlled comparison ($N=7$ edges, zero ablation). \css{} retains a small advantage over individual methods' top-7 edges, but the margin is substantially reduced from the uncontrolled comparison.}
\label{tab:ioi-size}
\end{table}

\begin{table}[t]
\centering
\small
\begin{tabular}{lccc}
\toprule
\textbf{Comparison} & \textbf{Orig.~$\Delta$} & \textbf{Size-ctrl $\Delta$} & \textbf{Retained} \\
\midrule
CSS vs ACDC & $0.081$ & $0.017$ & 42\% \\
CSS vs EAP-IG & $0.218$ & $0.031$ & 14\% \\
CSS vs EP & $0.350$ & $0.044$ & 13\% \\
\bottomrule
\end{tabular}
\caption{Size confound decomposition for IOI. Controlling for circuit size reduces the \css{} advantage by 58--87\%, indicating that the majority of the original advantage is attributable to circuit size rather than consensus quality.}
\label{tab:size-decomp}
\end{table}

After size control, \css{}$=1.0$ retains a residual advantage of $\Delta\cpr{}=0.017$ over the best single method (ACDC top-7). While directionally consistent with the consensus hypothesis, this margin is small and, as we show next, not statistically significant.

\paragraph{Bootstrap confidence intervals.}
Table~\ref{tab:bootstrap} presents bootstrap 95\% confidence intervals ($B=10{,}000$ resamples) for \cpr{} differences across all configurations. The size-controlled IOI advantage has a wide CI of $[-0.038, 0.071]$ ($p=0.27$), failing to exclude zero. Even the original (uncontrolled) IOI advantage is only marginally significant ($p=0.089$, CI $[-0.012, 0.174]$).

\begin{table}[t]
\centering
\small
\begin{tabular}{llcccc}
\toprule
\textbf{Task} & \textbf{Comparison} & \textbf{$\Delta$\cpr{}} & \textbf{95\% CI} & \textbf{$p$} & \textbf{Sig.} \\
\midrule
IOI & CSS vs ACDC (orig.) & $0.081$ & $[-0.01, 0.17]$ & $.089$ & No \\
IOI & CSS vs ACDC (size-ctrl) & $0.017$ & $[-0.04, 0.07]$ & $.270$ & No \\
IOI & CSS vs EAP-IG (orig.) & $0.218$ & $[0.10, 0.34]$ & $.002$ & Yes$^*$ \\
Arith. & CSS vs EAP-IG & $-0.066$ & $[-0.14, 0.01]$ & $.960$ & No \\
SVA & CSS vs ACDC & $0.013$ & $[-0.05, 0.08]$ & $.350$ & No \\
\bottomrule
\end{tabular}
\caption{Bootstrap 95\% CIs for \cpr{} differences ($B=10{,}000$). $^*$Significant but confounded by circuit size (see Table~\ref{tab:size-decomp}).}
\label{tab:bootstrap}
\end{table}

We therefore characterize IOI as a \emph{marginal} positive: \css{} shows a directionally correct but statistically weak advantage that is largely attributable to circuit size effects.

%-------------------------------------------------------------------
\subsection{Factual Recall: An Inconclusive Configuration}
\label{sec:factual}
%-------------------------------------------------------------------

Country-Capital factual recall on GPT-2 Small tests knowledge retrieval circuits, which are expected to be MLP-heavy. We apply the same $K=3$ methods with top-100 binarization. Table~\ref{tab:factual} summarizes the results.

\begin{table}[t]
\centering
\small
\begin{tabular}{lccc}
\toprule
\textbf{Circuit} & \textbf{Edges} & \textbf{\cpr{}$\uparrow$} & \textbf{\cmd{}$\downarrow$} \\
\midrule
\css{}$=1.0$ & 12 & $0.83$ & $0.41$ \\
\css{}$\geq 0.67$ & 38 & $0.79$ & $0.48$ \\
EAP-IG top-100 & 100 & $0.88$ & $0.32$ \\
EP top-100 & 100 & $0.82$ & $0.43$ \\
ACDC top-100 & 100 & $0.85$ & $0.37$ \\
\bottomrule
\end{tabular}
\caption{Factual Recall \cpr{} comparison (zero ablation). \css{}$=1.0$ does not outperform the best single method (EAP-IG, $\cpr{}=0.88$ vs $0.83$).}
\label{tab:factual}
\end{table}

The consensus circuit ($\css{}=1.0$, 12 edges) achieves $\cpr{}=0.83$, which is \emph{lower} than EAP-IG alone ($\cpr{}=0.88$, $\Delta=-0.05$). Pairwise Jaccard indices are low overall (EAP-IG$\times$EP: 0.42, EAP-IG$\times$ACDC: 0.18, EP$\times$ACDC: 0.15, three-way: 0.11), with 89\% of edges appearing in the off-diagonal region of the CSS distribution.

The failure pattern here differs from other configurations: methods disagree not because of binarization artifacts or scaling issues, but because of \emph{component-type preference}---gradient methods favor attention edges (where gradients through softmax are larger), while ACDC's ablation-based approach better localizes MLP contributions. This creates low Jaccard scores despite each method individually achieving reasonable \cpr{}. The result is inconclusive: \css{} neither helps nor catastrophically fails, but cannot recover the complementary information that different methods capture about different component types.

%-------------------------------------------------------------------
\subsection{Failure Mode 1: ACDC Binarization Pathology}
\label{sec:binarization}
%-------------------------------------------------------------------

ACDC produces circuits via iterative activation patching with an \texttt{isinf}-like threshold: edges whose ablation effect exceeds a threshold are retained. This creates a structural asymmetry with gradient-based methods: ACDC's output is already binary and typically much sparser.

Under the standard pipeline, \css{}$=1.0$ edges are necessarily a subset of ACDC's output (since ACDC contributes a binary $\{0,1\}$ vote). In configurations where ACDC selects fewer edges than the top-$N$ used for gradient methods, the consensus circuit degenerates to a subset of ACDC's circuit. We observed this pattern on IOI (5--6 of 7 consensus edges originate from ACDC), Arithmetic, and SVA.

Switching to uniform top-$N$ binarization across all methods breaks this degeneracy and shifts pairwise Jaccard by $+0.16$ to $+0.25$, confirming that the pathology is an artifact of binarization strategy rather than genuine method agreement.

%-------------------------------------------------------------------
\subsection{Failure Mode 2: Consensus Redundancy}
\label{sec:redundancy}
%-------------------------------------------------------------------

On Greater-Than Arithmetic and Subject-Verb Agreement (SVA), pairwise Jaccard indices are high (Arithmetic: max 0.695; SVA: max 0.653). When methods already agree substantially, consensus provides no new information---the intersection is essentially a single method's output with slightly fewer edges. For Arithmetic, \css{} produces $\Delta\cpr{}=-0.066$ versus the best single method (bootstrap CI $[-0.14, 0.01]$, $p=0.96$), indicating that the intersection actually \emph{removes} useful edges that not all methods select. SVA shows a near-zero effect ($\Delta=0.013$, CI $[-0.05, 0.08]$).

%-------------------------------------------------------------------
\subsection{Failure Mode 3: Scoring Granularity Mismatch}
\label{sec:granularity}
%-------------------------------------------------------------------

On Qwen-0.5B (${\sim}179$K edges), we observe a qualitatively different failure: near-zero inter-method overlap regardless of binarization. Root cause analysis reveals that both non-gradient methods degenerate to coarser-than-edge scoring:

\begin{table}[t]
\centering
\small
\begin{tabular}{lccc}
\toprule
\textbf{Method} & \textbf{Scoring Level} & \textbf{Eff.~Resolution} & \textbf{Unique / 100} \\
\midrule
EAP-IG & edge & 1 edge & 100 \\
EP & edge & 1 edge & 100 \\
IFR & block (layer) & ${\sim}200$ edges & ${\sim}3$--5 \\
ACDC & node (head/MLP) & ${\sim}500$--2K edges & ${\sim}8$--12 \\
\bottomrule
\end{tabular}
\caption{Scoring granularity comparison on Qwen-0.5B. Non-gradient methods assign identical scores to all edges within a component, producing block-structured selections with few unique score values.}
\label{tab:granularity}
\end{table}

IFR assigns identical scores to all edges within a layer block (${\sim}200$ edges per block), while ACDC assigns identical scores to all edges from the same attention head or MLP block (${\sim}500$--2{,}000 edges per node). Table~\ref{tab:granularity-jaccard} shows the resulting pairwise Jaccard values.

\begin{table}[t]
\centering
\small
\begin{tabular}{lcc}
\toprule
\textbf{Pair} & \textbf{Jaccard} & \textbf{Same Gran.?} \\
\midrule
EAP-IG $\times$ EP & 0.54 & Yes (edge--edge) \\
EAP-IG $\times$ IFR & 0.00 & No (edge--block) \\
EAP-IG $\times$ ACDC & 0.04 & No (edge--node) \\
EP $\times$ IFR & 0.00 & No (edge--block) \\
EP $\times$ ACDC & 0.04 & No (edge--node) \\
IFR $\times$ ACDC & 0.00 & No (block--node) \\
\bottomrule
\end{tabular}
\caption{Pairwise Jaccard on Qwen-0.5B by scoring granularity. Methods operating at different granularities have near-zero overlap.}
\label{tab:granularity-jaccard}
\end{table}

This is not an implementation bug but an architectural limitation: IFR's forward intervention is inherently layer-granular, and ACDC's activation patching operates at node granularity. With $K=3$ methods where the two gradient methods are highly correlated (Jaccard $=0.54$), the effective number of independent information sources drops from $K=3$ to $K \approx 1.5$, rendering consensus meaningless.

%-------------------------------------------------------------------
\subsection{Jaccard Diagnostic}
\label{sec:diagnostic}
%-------------------------------------------------------------------

Aggregating across all 5 task$\times$model configurations, we observe a negative relationship between maximum pairwise Jaccard and \css{} advantage ($\Delta\cpr{}$), as shown in Table~\ref{tab:jaccard-scatter}.

\begin{table}[t]
\centering
\small
\begin{tabular}{lccc}
\toprule
\textbf{Configuration} & \textbf{Max Jaccard} & \textbf{$\Delta$\cpr{}} & \textbf{CSS Helps?} \\
\midrule
IOI / GPT-2 & 0.307 & $+0.081$ & Yes (marginal) \\
Factual Recall / GPT-2 & 0.420 & $-0.050$ & No \\
IOI / Qwen-0.5B & 0.540 & $0.000$ & No \\
SVA / GPT-2 & 0.653 & $+0.013$ & No \\
Arithmetic / GPT-2 & 0.695 & $-0.066$ & No \\
\bottomrule
\end{tabular}
\caption{Max pairwise Jaccard vs \css{} advantage across all configurations. \css{} shows a positive effect only in the lowest-Jaccard configuration.}
\label{tab:jaccard-scatter}
\end{table}

A simple threshold rule emerges: configurations with max Jaccard $< 0.35$ benefit from \css{} (1/1), while those with max Jaccard $> 0.50$ never do (0/3). The gap zone ($0.35$--$0.50$, Factual Recall) yields an incorrect prediction. Spearman $\rho = -0.70$, but with $n=5$, $p=0.19$---far from significant.

We frame this as a \emph{motivated conjecture}: the pattern is directionally consistent and mechanistically plausible (low Jaccard implies diverse methods, a prerequisite for useful consensus), but $n=5$ configurations are insufficient for statistical validation. Practitioners may find the Jaccard check useful as a cheap pre-screening heuristic, but it should not be treated as a validated decision boundary.

%-------------------------------------------------------------------
\subsection{CIRCUS Comparison}
\label{sec:circus}
%-------------------------------------------------------------------

As an ablation, we compare \css{} (cross-method consensus) against CIRCUS-style within-method threshold diversity: varying the binarization threshold within a single method to generate multiple circuit ``voters.'' CIRCUS produces negative \cpr{} on all three tested tasks (IOI, Arithmetic, SVA), confirming that cross-method diversity provides signal that within-method threshold variation cannot replicate.

%===================================================================
\section{Discussion}
%===================================================================

\paragraph{Failure mode taxonomy.}
Our three failure modes map onto known conditions from ensemble theory \cite{kuncheva2003measures}. Mode~1 (binarization pathology) violates the \emph{independence} assumption: one method's binarization dominates the consensus. Mode~2 (high agreement) violates the \emph{diversity} requirement: highly correlated methods provide redundant votes. Mode~3 (granularity mismatch) violates the \emph{commensurability} assumption: methods must operate on the same feature space for intersection to be meaningful.

\paragraph{Practitioner guidelines.}
Based on our findings, we recommend: (1)~normalize binarization across methods (use top-$N$ uniformly rather than mixing native thresholds); (2)~compute pairwise Jaccard before investing in consensus---if max Jaccard $> 0.5$, consensus is unlikely to help; (3)~verify that all methods score at the same granularity, especially on models with $> 100$K edges.

%===================================================================
\section{Limitations}
\label{sec:limitations}
%===================================================================

\paragraph{Size confound.}
The sole positive configuration (IOI on GPT-2) is substantially confounded by circuit size. In the original comparison, \css{}$=1.0$ (7 edges) is evaluated against single-method top-100 circuits, creating a $14\times$ size disparity. After controlling for size (Table~\ref{tab:ioi-size}), the \css{} advantage shrinks from $\Delta=0.081$ to $\Delta=0.017$---a 79\% reduction. The residual advantage, while directionally correct, is not statistically significant ($p=0.27$). We cannot rule out that the entire \css{} effect on IOI is a size artifact.

\paragraph{Statistical power.}
Bootstrap confidence intervals are wide throughout: the IOI original-comparison CI spans $[-0.012, 0.174]$, with the lower bound near zero. Our evaluation set (100 prompts $\times$ 7 edges) provides limited statistical power. The Jaccard diagnostic analysis covers only $n=5$ configurations, insufficient for correlation testing ($p=0.19$). All quantitative claims should be interpreted as directional evidence, not definitive conclusions.

\paragraph{Task coverage.}
We evaluate on four tasks (IOI, Arithmetic, SVA, Factual Recall) across two models, totaling 5 complete configurations. This is insufficient to generalize failure mode prevalence or validate the Jaccard threshold. Notably, we have only one configuration with max Jaccard $< 0.35$, making it impossible to assess \css{}'s success rate in the low-Jaccard regime.

\paragraph{Ablation method.}
All evaluations use zero ablation. Zero ablation has known biases for small circuits: ablating $N-k$ edges to zero (where $k \ll N$) can amplify or suppress performance differences relative to mean ablation. We did not complete a systematic comparison of zero vs.\ mean ablation due to time constraints. Our results may be sensitive to this choice, particularly the IOI positive finding where the circuit contains only 7 edges.

\paragraph{Method selection.}
Our $K=3$ method set includes two gradient-based methods (EAP-IG, EP) that are inherently correlated (Jaccard $= 0.54$ on Qwen-0.5B). This reduces effective diversity. Alternative method sets (e.g., replacing EP with a non-gradient method) might yield different consensus behavior. We chose these methods for their public availability and widespread use, but acknowledge this limits the generality of our negative findings.

%===================================================================
\section{Conclusion}
%===================================================================

We presented a systematic evaluation of cross-method binary consensus (\css{}) for circuit discovery across four tasks and two model scales. Our main finding is negative: consensus succeeds marginally on only 1 of 5 configurations, with the advantage largely attributable to circuit size effects. We identified three mechanistically distinct failure modes---binarization pathology, consensus redundancy, and scoring granularity mismatch---each with a diagnostic signature. While our task and model coverage is limited ($n=5$ configurations) and statistical power is low, the failure modes are mechanistically grounded and yield concrete practitioner guidelines.

The broader lesson is that multi-method consensus is not a free lunch in circuit discovery. The conditions under which ensemble approaches improve over individual methods are specific and verifiable. We hope that making these failure modes explicit encourages the community to develop consensus methods that are robust to binarization artifacts, adaptive to method diversity, and aware of scoring granularity---or to invest in improving individual methods rather than ensembling flawed ones.

\paragraph{Future work.}
Validating the Jaccard diagnostic on a broader set of tasks and models is a priority. A systematic comparison of ablation methods (zero vs.\ mean) on all configurations would clarify whether the IOI positive finding is robust. Continuous consensus methods \cite{mondorf2025cross} that avoid binarization entirely may circumvent Failure Mode~1 and deserve comparative evaluation.

\begin{thebibliography}{20}

\bibitem[Bhaskar et~al.(2024)]{bhaskar2024finding}
Adithya Bhaskar, Alexander Wettig, Dan Friedman, and Danqi Chen.
\newblock Finding transformer circuits with edge pruning.
\newblock \emph{arXiv preprint arXiv:2406.16778}, 2024.

\bibitem[Chen et~al.(2026)]{chen2026pie}
Xiao Chen et~al.
\newblock PIE: Pruning-based iterative ensembles for circuit discovery.
\newblock \emph{arXiv preprint arXiv:2604.16889}, 2026.

\bibitem[Conmy et~al.(2023)]{conmy2023towards}
Arthur Conmy, Augustine Mavor-Parker, Aaquib Lynch, Stefan Heimersheim, and Adri\`{a} Garriga-Alonso.
\newblock Towards automated circuit discovery for mechanistic interpretability.
\newblock In \emph{NeurIPS}, 2023.

\bibitem[Kuncheva and Whitaker(2003)]{kuncheva2003measures}
Ludmila~I. Kuncheva and Christopher~J. Whitaker.
\newblock Measures of diversity in classifier ensembles and their relationship with the ensemble accuracy.
\newblock \emph{Machine Learning}, 51(2):181--207, 2003.

\bibitem[Mondorf et~al.(2025)]{mondorf2025cross}
Manuel Mondorf et~al.
\newblock Cross-method score averaging for circuit discovery.
\newblock \emph{arXiv preprint arXiv:2510.06811}, 2025.

\bibitem[Nogueira and Brown(2016)]{nogueira2016measuring}
Sarah Nogueira and Gavin Brown.
\newblock Measuring the stability of feature selection.
\newblock \emph{JMLR}, 2016.

\bibitem[Parekh et~al.(2026)]{parekh2026circus}
Jayneel Parekh et~al.
\newblock CIRCUS: Within-method threshold diversity for circuit discovery.
\newblock \emph{arXiv preprint}, 2026.

\bibitem[Syed et~al.(2023)]{syed2023attribution}
Aaquib Syed, Can Rager, and Arthur Conmy.
\newblock Attribution patching outperforms automated circuit discovery.
\newblock \emph{arXiv preprint arXiv:2310.10348}, 2023.

\bibitem[Wang et~al.(2023)]{wang2023interpretability}
Kevin Wang, Alexandre Variengien, Arthur Conmy, Buck Shlegeris, and Jacob Steinhardt.
\newblock Interpretability in the wild: a circuit for indirect object identification in GPT-2 Small.
\newblock In \emph{ICLR}, 2023.

\end{thebibliography}

\appendix
\section{Experimental Details}
\label{app:details}

All experiments use the TransformerLens library for model hooks and circuit extraction. For EAP-IG, we use 10-step integrated gradients. For Edge Pruning, we train continuous masks for 500 steps with $L_1$ regularization ($\lambda=0.01$). For ACDC, we use the default threshold ($\tau=0.01$) with iterative node-level patching.

Evaluation uses 100 prompts per task, sampled from the standard evaluation sets. Bootstrap confidence intervals use $B=10{,}000$ resamples with bias-corrected and accelerated (BCa) percentile intervals.

GPT-2 Small has 12 layers, 12 attention heads, and approximately 32K edges in the computational graph. Qwen-0.5B has 24 layers, 16 attention heads, and approximately 179K edges.

\end{document}
